\documentclass[11pt]{article}

\usepackage[margin=1in]{geometry}
\usepackage[T1]{fontenc}
\usepackage[utf8]{inputenc}
\usepackage{lmodern}
\usepackage{enumitem}
\usepackage{hyperref}

\hypersetup{
  colorlinks=true,
  linkcolor=blue!50!black,
  urlcolor=blue!50!black,
  citecolor=blue!50!black,
  pdftitle={Discovery Breakthrough Profile for p-vs-np}
}

\title{Discovery Breakthrough Profile for \texttt{p-vs-np}}
\author{Neuro-symbolic AI Group\\Idiap Research Institute}
\date{March 18, 2026}

\begin{document}
\maketitle

\begin{abstract}
This report classifies the blocked discovery frontier by the conceptual breakthrough required to make progress. It reframes the frontier from missing theorems to breakthrough families such as realization, transport, alignment, explicitness, constructive instantiation, and semantic bridging.
\end{abstract}

\section{Summary}
\begin{itemize}[leftmargin=*,itemsep=0.2em]
\item The discovery layer should stop restating blocked lemmas as if they were already search strategies.
\item Most frontier pressure is now in bridge-layer breakthroughs: realization, transport, alignment, explicitness, and semantic bridging.
\item Target-grounded routes should be treated as activation-theorem searches, not as target-theorem discovery.
\item The highest-pressure breakthrough family is 'abstraction-search' via obstruction class 'mcsp-magnification'.
\end{itemize}

\section{mcsp-magnification}
\textbf{Breakthrough family.} abstraction-search

\textbf{Pressure score.} 5.818

\textbf{Required move.} Find a sharper abstraction that turns the blocked route into a typed, inspectable theorem target.

\textbf{Abstraction required.} A theorem-family-specific carrier, invariant, or bridge law that adds new mathematical content.

\textbf{Minimal theorem shape.} State a route-specific theorem that would resolve the blocked interface and make the next proof obligation materially sharper.

\textbf{Major subresult type.} exploratory-conjecture

\subsection*{Search Modes}
\begin{itemize}[leftmargin=*,itemsep=0.2em]
\item typed abstraction search
\item bridge split
\item counterfactual sharpening
\end{itemize}
\subsection*{Anti-Patterns}
\begin{itemize}[leftmargin=*,itemsep=0.2em]
\item repeating the same blocked statement without new typed structure
\end{itemize}
\section{carrier-realization}
\textbf{Breakthrough family.} realization

\textbf{Pressure score.} 5.818

\textbf{Required move.} Turn a typed but schematic carrier into an explicit mathematical object family with admissibility conditions and a forcing theorem.

\textbf{Abstraction required.} An explicit obstruction family, witness family, surrogate promise problem, or quantitative measure that is realized from the source theorem rather than only named by role.

\textbf{Minimal theorem shape.} If source theorem S holds in regime R, then one can construct an explicit carrier C with admissibility conditions A such that C preserves invariant I and supports a route-specific forcing lemma to target T.

\textbf{Major subresult type.} bridge object/tool

\subsection*{Search Modes}
\begin{itemize}[leftmargin=*,itemsep=0.2em]
\item explicit carrier construction
\item equivalence search between live carrier encodings
\item surrogate problem design with forcing semantics
\end{itemize}
\subsection*{Anti-Patterns}
\begin{itemize}[leftmargin=*,itemsep=0.2em]
\item renaming the carrier without adding realizability conditions
\item adding more obligations without exhibiting a concrete carrier family
\end{itemize}
\section{formula-size-transport}
\textbf{Breakthrough family.} transport

\textbf{Pressure score.} 5.060

\textbf{Required move.} Find a theorem that preserves the relevant invariant across a cross-model move into formula-size.

\textbf{Abstraction required.} A transport law, compilation theorem, or complexity measure that is stable under the source-to-target model shift.

\textbf{Minimal theorem shape.} For every admissible carrier C in the nonuniform-circuit regime, there is a transformed carrier C' in the formula-size regime preserving invariant I and retaining enough quantitative strength to imply target T.

\textbf{Major subresult type.} transfer principle

\subsection*{Search Modes}
\begin{itemize}[leftmargin=*,itemsep=0.2em]
\item cross-model transport principle
\item invariant-preserving compilation theorem
\item representation normal form for formula-size transfer
\end{itemize}
\subsection*{Anti-Patterns}
\begin{itemize}[leftmargin=*,itemsep=0.2em]
\item assuming the target theorem already provides the transport
\item strengthening the source lower bound without addressing the model shift
\end{itemize}
\section{target-grounding}
\textbf{Breakthrough family.} bridge-specificity

\textbf{Pressure score.} 5.014

\textbf{Required move.} Stop treating the target theorem as the open problem and relocate search to the bridge that would activate it.

\textbf{Abstraction required.} No new target theorem; discovery should instead seek the carrier, transfer law, or alignment theorem that makes the target theorem applicable.

\textbf{Minimal theorem shape.} Target theorem T is fixed. The real theorem under search is the local activation theorem A that verifies T's hypotheses from current artifacts.

\textbf{Major subresult type.} research recommendation

\subsection*{Search Modes}
\begin{itemize}[leftmargin=*,itemsep=0.2em]
\item activation theorem search
\item hypothesis verification search
\item bridge-localization refinement
\end{itemize}
\subsection*{Anti-Patterns}
\begin{itemize}[leftmargin=*,itemsep=0.2em]
\item rediscovering the target theorem
\item claiming novelty at the target instead of at the activation bridge
\end{itemize}
\section{local-reduction-barrier}
\textbf{Breakthrough family.} barrier-aware transfer

\textbf{Pressure score.} 1.984

\textbf{Required move.} Design a transfer mechanism that survives known local-reduction barriers, or prove that no such mechanism exists under explicit assumptions.

\textbf{Abstraction required.} A barrier-aware reduction gadget family or an impossibility theorem explaining why the reduction regime cannot be crossed.

\textbf{Minimal theorem shape.} Either construct a local reduction family G satisfying locality bounds L and preserving witness structure W, or prove that any such G would contradict barrier theorem B.

\textbf{Major subresult type.} barrier result

\subsection*{Search Modes}
\begin{itemize}[leftmargin=*,itemsep=0.2em]
\item barrier-aware gadget design
\item impossibility theorem search
\item reduction regime decomposition
\end{itemize}
\subsection*{Anti-Patterns}
\begin{itemize}[leftmargin=*,itemsep=0.2em]
\item ignoring the known local-reduction barrier theorem
\item using generic NP-hardness language without locality constraints
\end{itemize}
\section{target-alignment}
\textbf{Breakthrough family.} alignment

\textbf{Pressure score.} 1.897

\textbf{Required move.} Identify or prove an equivalence between the current carrier class and the object class required by the target theorem.

\textbf{Abstraction required.} An alignment theorem, correspondence theorem, or mediating normal form linking the local carrier to the target's sparse-problem or distinguisher objects.

\textbf{Minimal theorem shape.} Every admissible carrier C from the local route induces an object D in the target theorem's class, and D preserves the quantitative property needed to activate the target theorem.

\textbf{Major subresult type.} transfer principle

\subsection*{Search Modes}
\begin{itemize}[leftmargin=*,itemsep=0.2em]
\item object-class correspondence search
\item normal-form mediation
\item dual characterization of the target object
\end{itemize}
\subsection*{Anti-Patterns}
\begin{itemize}[leftmargin=*,itemsep=0.2em]
\item treating topical similarity as object equivalence
\item adding a stronger carrier without specifying the target object class map
\end{itemize}
\section{general-magnification}
\textbf{Breakthrough family.} alignment

\textbf{Pressure score.} 1.897

\textbf{Required move.} Map a problem-specific carrier into a general magnification framework without losing the quantitative threshold that matters.

\textbf{Abstraction required.} A general-purpose sparse-problem or distinguisher formulation that faithfully receives the local carrier.

\textbf{Minimal theorem shape.} Problem-specific object C induces general object G with preserved sparsity or approximation parameters sufficient to trigger the general magnification theorem.

\textbf{Major subresult type.} transfer principle

\subsection*{Search Modes}
\begin{itemize}[leftmargin=*,itemsep=0.2em]
\item generalization by faithful embedding
\item parameter-preserving reformulation
\item problem-specific to generic theorem bridge
\end{itemize}
\subsection*{Anti-Patterns}
\begin{itemize}[leftmargin=*,itemsep=0.2em]
\item moving to the generic theorem without a preserved parameter map
\item using the general theorem only as motivation rather than as a typed target
\end{itemize}
\section{pit-upgrade}
\textbf{Breakthrough family.} explicitness

\textbf{Pressure score.} 1.752

\textbf{Required move.} Upgrade a semantic or existential algebraic object into an explicit algorithmic derandomization premise.

\textbf{Abstraction required.} An explicit generator, certificate family, or algebraic witness system that is constructive enough for PIT rather than merely semantic.

\textbf{Minimal theorem shape.} If semantic object C exists for the relevant circuit class, then one can uniformly and efficiently extract an explicit structure G that yields deterministic PIT for that class.

\textbf{Major subresult type.} bridge object/tool

\subsection*{Search Modes}
\begin{itemize}[leftmargin=*,itemsep=0.2em]
\item explicit extractor or generator search
\item certificate map construction
\item semantic-to-algorithmic upgrade theorem
\end{itemize}
\subsection*{Anti-Patterns}
\begin{itemize}[leftmargin=*,itemsep=0.2em]
\item equating a black-box hitting-set statement with an explicit generator
\item rephrasing derandomization as semantics without uniformity or efficiency
\end{itemize}
\section{schema-to-examples}
\textbf{Breakthrough family.} constructive instantiation

\textbf{Pressure score.} 1.641

\textbf{Required move.} Turn an abstract schema into explicit natural example families or prove that the schema cannot instantiate naturally in the current way.

\textbf{Abstraction required.} A construction template, example-generating reduction family, or a canonicalization principle for natural examples.

\textbf{Minimal theorem shape.} From abstract schema S, construct an explicit family of natural instances E\_n satisfying properties P and realizing target theorem T.

\textbf{Major subresult type.} standalone theorem

\subsection*{Search Modes}
\begin{itemize}[leftmargin=*,itemsep=0.2em]
\item canonical example construction
\item template instantiation
\item example-preserving reduction search
\end{itemize}
\subsection*{Anti-Patterns}
\begin{itemize}[leftmargin=*,itemsep=0.2em]
\item restating the schema as if it already supplied examples
\item listing examples without a constructive route from the local abstraction
\end{itemize}
\section{characterization-bridge}
\textbf{Breakthrough family.} semantic bridge

\textbf{Pressure score.} 1.621

\textbf{Required move.} Bridge a local definitional object to a global characterization theorem.

\textbf{Abstraction required.} A semantic correspondence, duality, witness extraction principle, or completeness notion that turns the local definition into the target biconditional.

\textbf{Minimal theorem shape.} Local definitional object D has property P if and only if global characterization condition Q holds, with both directions proved at the right level of abstraction.

\textbf{Major subresult type.} transfer principle

\subsection*{Search Modes}
\begin{itemize}[leftmargin=*,itemsep=0.2em]
\item semantic correspondence theorem
\item duality or normal form search
\item two-way implication decomposition
\end{itemize}
\subsection*{Anti-Patterns}
\begin{itemize}[leftmargin=*,itemsep=0.2em]
\item normalizing syntax without adding semantic force
\item treating a definition as if it already contained the characterization
\end{itemize}
\section{explicitness-gap}
\textbf{Breakthrough family.} explicitness

\textbf{Pressure score.} 1.342

\textbf{Required move.} Add uniform, explicit, efficiently constructible content to an otherwise semantic object.

\textbf{Abstraction required.} A uniform construction rule, extractor, or explicit family indexed by input length and class parameters.

\textbf{Minimal theorem shape.} For every length parameter n, the existential object C\_n can be constructed or certified explicitly in time class X while preserving property P.

\textbf{Major subresult type.} bridge object/tool

\subsection*{Search Modes}
\begin{itemize}[leftmargin=*,itemsep=0.2em]
\item uniformization
\item effective construction
\item certificate extraction
\end{itemize}
\subsection*{Anti-Patterns}
\begin{itemize}[leftmargin=*,itemsep=0.2em]
\item accepting existence without algorithmic access
\item shifting to a broader semantic object that is still non-explicit
\end{itemize}

\end{document}
